% !PRINTED 231
\weberpage{231}

%% [running head, not body text -- commented out 2026-08-07] Page furniture (running
%% head, folio, signature) is excluded from this edition by standing policy.
%% \S.\,74. \hfill Principal equation of the fifth degree \hfill 231

%<<CONT
\weq{2}{
\begin{array}{rcl@{\qquad}rcl}
5\,B_{0,0} &=& 4\,a_4^2 - 10\,a_3\,a_5, & 5\,B_{0,1} &=& 3\,a_3\,a_4 \\[1ex]
5\,B_{1,1} &=& 6\,a_3^2, & 5\,B_{0,2} &=& 0 \\[1ex]
5\,B_{2,2} &=& -20\,a_0\,a_4, & 5\,B_{0,3} &=& -25\,a_0\,a_5 \\[1ex]
5\,B_{3,3} &=& 0 & 5\,B_{1,2} &=& -25\,a_0\,a_5 \\[1ex]
5\,B_{1,3} &=& -20\,a_0\,a_4, & 5\,B_{2,3} &=& -15\,a_0\,a_3.
\end{array}
}
If one computes the determinant from these, one obtains
without difficulty, by \S.\,73, (12), for the discriminant
of the principal equation of the fifth degree
\[
\text{(3)} \qquad D = a_0^2\bigl(5^5\,a_0^2\,a_5^4 + 2^8\,a_0\,a_4^5 + 2250\,a_0\,a_3^2\,a_4\,a_5^2
- 1600\,a_0\,a_3\,a_4^3\,a_5 + 108\,a_3^5\,a_5 - 27\,a_3^4\,a_4^2\bigr)\,.
\]
Now since $B_{3,3} = 0$ here, we have a system of values
of the $t$ for which the Bezoutiant vanishes, namely
\[
t_0 = 0, \quad t_1 = 0, \quad t_2 = 0\,.
\]
We now take our Tschirnhausen substitution in the following form [\S.\,68, (12)]
\[
\text{(4)} \qquad y = t_3\,F_0 + t_2\,(\alpha_2\,F_1 + \alpha_1\,F_2 + \alpha_0\,F_3)
\]
and we wish to dispose of the $\alpha_0, \alpha_1, \alpha_2$ in such a way that the equation
of the fifth degree for $y$ becomes a principal equation independently of $t_3, t_2$ --- that is, that $S(y^2)$ vanishes identically for all $t_2, t_3$.

\indent Since $S(F_0^2) = B_{3,3}$ vanishes here, this
requirement is satisfied, by \S.\,73, (11), if we determine $\alpha_0, \alpha_1, \alpha_2$ so that
\[
\text{(5)} \qquad \alpha_0^2\,B_{0,0} + \alpha_1^2\,B_{1,1} + \alpha_2^2\,B_{2,2}
+ 2\,\alpha_0\,\alpha_1\,B_{0,1} + 2\,\alpha_0\,\alpha_2\,B_{0,2} + 2\,\alpha_1\,\alpha_2\,B_{1,2} = 0\,,
\]
\[
\text{(6)} \qquad \alpha_0\,B_{0,3} + \alpha_1\,B_{1,3} + \alpha_2\,B_{2,3} = 0\,.
\]
These equations lead, by elimination of one of the
three quantities $\alpha$, to a quadratic equation for the ratio of the other two. In itself it is not necessary
to exclude any particular case, not even the one in which
$B_{0,3}, B_{1,3}, B_{2,3}$ all three vanish; one would then merely have
to take for the $\alpha$ some solution of equation (5).

\indent But in order not to be long-winded, we shall assume
that $B_{0,3}$ is different from zero, because the cases thereby excluded, in which a root of the given equation becomes zero or infinite, really offer no interest
here, since they come back to an equation of lower degree. The treatment, moreover, remains entirely the same if
%CONT>>
